%% ECON 282E -- Session 5, Deck B5: Determining the Coefficients
%% Source: Quant_Macro Lec.9 frames 96-108 (L2415-2742).  Follows
%% Fernandez-Villaverde & Guerron-Quintana, Solution Methods IV pp.39-47;
%% rewritten here and cited.  Source defects are recorded in slides/SOURCES.md,
%% not on the slides.
%% Numbers from tools/figures/s05_numbers.json, key "collocation".

%% ECON 282E -- Foundations of Macroeconomics (UC Riverside, Fall 2026)
%% Shared Beamer preamble for every deck in slides/S01 ... slides/S10.
%%
%% Usage (a deck lives in slides/S0X/ and is compiled from that folder):
%%   %% ECON 282E -- Foundations of Macroeconomics (UC Riverside, Fall 2026)
%% Shared Beamer preamble for every deck in slides/S01 ... slides/S10.
%%
%% Usage (a deck lives in slides/S0X/ and is compiled from that folder):
%%   %% ECON 282E -- Foundations of Macroeconomics (UC Riverside, Fall 2026)
%% Shared Beamer preamble for every deck in slides/S01 ... slides/S10.
%%
%% Usage (a deck lives in slides/S0X/ and is compiled from that folder):
%%   \input{../preamble/course_preamble.tex}
%%   \decktitle{1}{A1}{Who I Am \& Why This Course}{September 24, 2026}
%%   \begin{document}
%%   \makedecktitle
%%   ...
%%
%% Adapted from Courses/AI-ECON-2026/source/shared_preamble.tex (Metropolis theme,
%% concept/caution/example/takeaway boxes, \sectiondivider). Changes for this course:
%%   * course title block (\decktitle / \makedecktitle) instead of \makebeamertitle
%%   * \zh{...} typesets Chinese (book titles) under pdflatex via CJKutf8
%%   * researchbox for the "Research in the wild" frames
%%   * section pages off; TikZ libraries and the course map loaded here
%% Build with pdflatex only (two passes); tools/build_slides.py does this for you.

\documentclass[10pt,english,aspectratio=169]{beamer}

\usepackage[T1]{fontenc}
\usepackage{textcomp}
\usepackage[utf8]{inputenc}
\usepackage{babel}
\usepackage{CJKutf8}

\usepackage{amsmath, amssymb, amsthm, graphicx}
\usepackage{booktabs, multirow, tabularx}
\usepackage{tikz}
\usetikzlibrary{positioning,arrows.meta,shapes,calc,fit,backgrounds}
\usepackage{pgfplots}
\pgfplotsset{compat=1.16}
\usepackage[most]{tcolorbox}
\tcbuselibrary{skins, breakable}
\usepackage{listings}
\usepackage{xcolor}

\lstset{
  basicstyle=\ttfamily\small,
  keywordstyle=\color{blue!80!black}\bfseries,
  commentstyle=\color{green!50!black}\itshape,
  stringstyle=\color{orange!80!black},
  backgroundcolor=\color{gray!8},
  frame=single,
  rulecolor=\color{gray!40},
  breaklines=true,
  showstringspaces=false,
  numbers=none,
}

\usetheme{metropolis}
\usefonttheme{professionalfonts}
\metroset{sectionpage=none}

\definecolor{LightGray}{RGB}{230,230,230}
\definecolor{RedTitle}{RGB}{0,0,200}
\definecolor{VeryLightGray}{RGB}{245,245,245}
\definecolor{DarkText}{RGB}{0,0,0}
\definecolor{UCRBlue}{RGB}{0,45,106}
\definecolor{UCRGold}{RGB}{241,171,0}

% Stage colours of the course arc (used by the course map and elsewhere)
\colorlet{StageTrunk}{blue!12}
\colorlet{StageBranch}{cyan!14}
\colorlet{StageMachine}{gray!18}
\colorlet{StageWall}{red!14}
\colorlet{StageTools}{orange!20}
\colorlet{StageData}{green!16}
\colorlet{StageAgents}{violet!14}

\hypersetup{bookmarks=false,breaklinks=true,pdfborder={0 0 0},colorlinks=true,
  linkcolor=DarkText,urlcolor=RedTitle}

\setbeamercolor{background canvas}{bg=white}
\setbeamercolor{title}{fg=RedTitle}
\setbeamercolor{frametitle}{bg=VeryLightGray,fg=DarkText}
\setbeamercolor{normal text}{fg=DarkText}
\setbeamercolor{alerted text}{fg=RedTitle}
\setbeamersize{text margin left=5mm, text margin right=5mm}
\addtobeamertemplate{frametitle}{}{\vspace{-0.5em}}

\setbeamertemplate{navigation symbols}{}
\setbeamertemplate{footline}{%
  \hspace{5mm}{\usebeamerfont{page number in head/foot}\color{black!45}%
  ECON 282E \,$\cdot$\, \insertshortdeck}\hfill%
  \usebeamercolor[fg]{page number in head/foot}%
  \usebeamerfont{page number in head/foot}%
  \insertframenumber\,/\,\inserttotalframenumber\kern1em\vskip2pt%
}

% ---------------------------------------------------------------------------
% Chinese text under pdflatex (book titles etc.)
\newcommand{\zh}[1]{\begin{CJK*}{UTF8}{gbsn}#1\end{CJK*}}

% ---------------------------------------------------------------------------
% Deck title block
%   \decktitle{session}{deck id}{title}{date}
\newcommand{\insertshortdeck}{}
\newcommand{\decksession}{}
\newcommand{\deckid}{}
\newcommand{\decktitletext}{}
\newcommand{\deckdate}{}
\newcommand{\decktitle}[4]{%
  \renewcommand{\decksession}{#1}%
  \renewcommand{\deckid}{#2}%
  \renewcommand{\decktitletext}{#3}%
  \renewcommand{\deckdate}{#4}%
  \renewcommand{\insertshortdeck}{Session #1 \,$\cdot$\, Deck #1.#2}%
  \title{#3}%
  \author{Zhigang Feng}%
  \date{#4}%
}
\newcommand{\makedecktitle}{%
  \begin{frame}[plain,noframenumbering]
    \begin{tikzpicture}[remember picture,overlay]
      \fill[UCRBlue] (current page.north west) rectangle ([yshift=-0.55cm]current page.north east);
      \fill[UCRGold] ([yshift=-0.55cm]current page.north west) rectangle ([yshift=-0.62cm]current page.north east);
      \node[anchor=west, text=white, font=\small\bfseries] at ([xshift=5mm,yshift=-0.275cm]current page.north west)
        {ECON 282E \,$\cdot$\, Foundations of Macroeconomics};
      \node[anchor=east, text=white, font=\small] at ([xshift=-5mm,yshift=-0.275cm]current page.north east)
        {UC Riverside \,$\cdot$\, Fall 2026};
    \end{tikzpicture}
    \vspace*{1.1cm}
    {\usebeamerfont{title}\color{black!55}\large Session \decksession\ \,$\cdot$\, Deck \decksession.\deckid\par}
    \vspace{0.25cm}
    {\usebeamerfont{title}\color{RedTitle}\LARGE\bfseries \decktitletext\par}
    \vspace{0.7cm}
    {\large Zhigang Feng\par}
    \vspace{0.15cm}
    {\color{black!65}\deckdate\par}
    \vfill
    {\footnotesize\itshape\color{black!55}Build it on paper $\;\to\;$ solve it on the machine $\;\to\;$ push it to the frontier with AI\par}
  \end{frame}}

% ---------------------------------------------------------------------------
% Section divider (plain frame, not counted)
\newcommand{\sectiondivider}[2]{%
\begin{frame}[plain,noframenumbering]
\begin{center}
\vfill
{\Large\textbf{#1}\par}
\vspace{0.35cm}
{\normalsize #2\par}
\vfill
\end{center}
\end{frame}
}

% ---------------------------------------------------------------------------
% Boxes
\newtcolorbox{conceptbox}[1][]{colback=blue!3, colframe=RedTitle!55, boxrule=0.35mm,
  arc=1mm, left=2mm, right=2mm, top=1mm, bottom=1mm, #1}
\newtcolorbox{cautionbox}[1][]{colback=red!3, colframe=red!50!black, boxrule=0.35mm,
  arc=1mm, left=2mm, right=2mm, top=1mm, bottom=1mm, #1}
\newtcolorbox{examplebox}[1][]{colback=green!3, colframe=green!45!black, boxrule=0.35mm,
  arc=1mm, left=2mm, right=2mm, top=1mm, bottom=1mm, #1}
\newtcolorbox{takeawaybox}[1][]{colback=VeryLightGray, colframe=RedTitle!55, boxrule=0.35mm,
  arc=1mm, left=2mm, right=2mm, top=1mm, bottom=1mm, #1}
\newtcolorbox{researchbox}[1][]{enhanced, colback=violet!4, colframe=violet!60!black,
  boxrule=0.35mm, arc=1mm, left=2mm, right=2mm, top=1mm, bottom=1mm,
  fonttitle=\bfseries\small, title={Research in the wild}, #1}

\newcommand{\compacttables}{\renewcommand{\arraystretch}{1.15}}
\newcommand{\src}[1]{{\tiny\color{black!55}#1}}

% ---------------------------------------------------------------------------
\input{../preamble/course_map.tex}

%%   \decktitle{1}{A1}{Who I Am \& Why This Course}{September 24, 2026}
%%   \begin{document}
%%   \makedecktitle
%%   ...
%%
%% Adapted from Courses/AI-ECON-2026/source/shared_preamble.tex (Metropolis theme,
%% concept/caution/example/takeaway boxes, \sectiondivider). Changes for this course:
%%   * course title block (\decktitle / \makedecktitle) instead of \makebeamertitle
%%   * \zh{...} typesets Chinese (book titles) under pdflatex via CJKutf8
%%   * researchbox for the "Research in the wild" frames
%%   * section pages off; TikZ libraries and the course map loaded here
%% Build with pdflatex only (two passes); tools/build_slides.py does this for you.

\documentclass[10pt,english,aspectratio=169]{beamer}

\usepackage[T1]{fontenc}
\usepackage{textcomp}
\usepackage[utf8]{inputenc}
\usepackage{babel}
\usepackage{CJKutf8}

\usepackage{amsmath, amssymb, amsthm, graphicx}
\usepackage{booktabs, multirow, tabularx}
\usepackage{tikz}
\usetikzlibrary{positioning,arrows.meta,shapes,calc,fit,backgrounds}
\usepackage{pgfplots}
\pgfplotsset{compat=1.16}
\usepackage[most]{tcolorbox}
\tcbuselibrary{skins, breakable}
\usepackage{listings}
\usepackage{xcolor}

\lstset{
  basicstyle=\ttfamily\small,
  keywordstyle=\color{blue!80!black}\bfseries,
  commentstyle=\color{green!50!black}\itshape,
  stringstyle=\color{orange!80!black},
  backgroundcolor=\color{gray!8},
  frame=single,
  rulecolor=\color{gray!40},
  breaklines=true,
  showstringspaces=false,
  numbers=none,
}

\usetheme{metropolis}
\usefonttheme{professionalfonts}
\metroset{sectionpage=none}

\definecolor{LightGray}{RGB}{230,230,230}
\definecolor{RedTitle}{RGB}{0,0,200}
\definecolor{VeryLightGray}{RGB}{245,245,245}
\definecolor{DarkText}{RGB}{0,0,0}
\definecolor{UCRBlue}{RGB}{0,45,106}
\definecolor{UCRGold}{RGB}{241,171,0}

% Stage colours of the course arc (used by the course map and elsewhere)
\colorlet{StageTrunk}{blue!12}
\colorlet{StageBranch}{cyan!14}
\colorlet{StageMachine}{gray!18}
\colorlet{StageWall}{red!14}
\colorlet{StageTools}{orange!20}
\colorlet{StageData}{green!16}
\colorlet{StageAgents}{violet!14}

\hypersetup{bookmarks=false,breaklinks=true,pdfborder={0 0 0},colorlinks=true,
  linkcolor=DarkText,urlcolor=RedTitle}

\setbeamercolor{background canvas}{bg=white}
\setbeamercolor{title}{fg=RedTitle}
\setbeamercolor{frametitle}{bg=VeryLightGray,fg=DarkText}
\setbeamercolor{normal text}{fg=DarkText}
\setbeamercolor{alerted text}{fg=RedTitle}
\setbeamersize{text margin left=5mm, text margin right=5mm}
\addtobeamertemplate{frametitle}{}{\vspace{-0.5em}}

\setbeamertemplate{navigation symbols}{}
\setbeamertemplate{footline}{%
  \hspace{5mm}{\usebeamerfont{page number in head/foot}\color{black!45}%
  ECON 282E \,$\cdot$\, \insertshortdeck}\hfill%
  \usebeamercolor[fg]{page number in head/foot}%
  \usebeamerfont{page number in head/foot}%
  \insertframenumber\,/\,\inserttotalframenumber\kern1em\vskip2pt%
}

% ---------------------------------------------------------------------------
% Chinese text under pdflatex (book titles etc.)
\newcommand{\zh}[1]{\begin{CJK*}{UTF8}{gbsn}#1\end{CJK*}}

% ---------------------------------------------------------------------------
% Deck title block
%   \decktitle{session}{deck id}{title}{date}
\newcommand{\insertshortdeck}{}
\newcommand{\decksession}{}
\newcommand{\deckid}{}
\newcommand{\decktitletext}{}
\newcommand{\deckdate}{}
\newcommand{\decktitle}[4]{%
  \renewcommand{\decksession}{#1}%
  \renewcommand{\deckid}{#2}%
  \renewcommand{\decktitletext}{#3}%
  \renewcommand{\deckdate}{#4}%
  \renewcommand{\insertshortdeck}{Session #1 \,$\cdot$\, Deck #1.#2}%
  \title{#3}%
  \author{Zhigang Feng}%
  \date{#4}%
}
\newcommand{\makedecktitle}{%
  \begin{frame}[plain,noframenumbering]
    \begin{tikzpicture}[remember picture,overlay]
      \fill[UCRBlue] (current page.north west) rectangle ([yshift=-0.55cm]current page.north east);
      \fill[UCRGold] ([yshift=-0.55cm]current page.north west) rectangle ([yshift=-0.62cm]current page.north east);
      \node[anchor=west, text=white, font=\small\bfseries] at ([xshift=5mm,yshift=-0.275cm]current page.north west)
        {ECON 282E \,$\cdot$\, Foundations of Macroeconomics};
      \node[anchor=east, text=white, font=\small] at ([xshift=-5mm,yshift=-0.275cm]current page.north east)
        {UC Riverside \,$\cdot$\, Fall 2026};
    \end{tikzpicture}
    \vspace*{1.1cm}
    {\usebeamerfont{title}\color{black!55}\large Session \decksession\ \,$\cdot$\, Deck \decksession.\deckid\par}
    \vspace{0.25cm}
    {\usebeamerfont{title}\color{RedTitle}\LARGE\bfseries \decktitletext\par}
    \vspace{0.7cm}
    {\large Zhigang Feng\par}
    \vspace{0.15cm}
    {\color{black!65}\deckdate\par}
    \vfill
    {\footnotesize\itshape\color{black!55}Build it on paper $\;\to\;$ solve it on the machine $\;\to\;$ push it to the frontier with AI\par}
  \end{frame}}

% ---------------------------------------------------------------------------
% Section divider (plain frame, not counted)
\newcommand{\sectiondivider}[2]{%
\begin{frame}[plain,noframenumbering]
\begin{center}
\vfill
{\Large\textbf{#1}\par}
\vspace{0.35cm}
{\normalsize #2\par}
\vfill
\end{center}
\end{frame}
}

% ---------------------------------------------------------------------------
% Boxes
\newtcolorbox{conceptbox}[1][]{colback=blue!3, colframe=RedTitle!55, boxrule=0.35mm,
  arc=1mm, left=2mm, right=2mm, top=1mm, bottom=1mm, #1}
\newtcolorbox{cautionbox}[1][]{colback=red!3, colframe=red!50!black, boxrule=0.35mm,
  arc=1mm, left=2mm, right=2mm, top=1mm, bottom=1mm, #1}
\newtcolorbox{examplebox}[1][]{colback=green!3, colframe=green!45!black, boxrule=0.35mm,
  arc=1mm, left=2mm, right=2mm, top=1mm, bottom=1mm, #1}
\newtcolorbox{takeawaybox}[1][]{colback=VeryLightGray, colframe=RedTitle!55, boxrule=0.35mm,
  arc=1mm, left=2mm, right=2mm, top=1mm, bottom=1mm, #1}
\newtcolorbox{researchbox}[1][]{enhanced, colback=violet!4, colframe=violet!60!black,
  boxrule=0.35mm, arc=1mm, left=2mm, right=2mm, top=1mm, bottom=1mm,
  fonttitle=\bfseries\small, title={Research in the wild}, #1}

\newcommand{\compacttables}{\renewcommand{\arraystretch}{1.15}}
\newcommand{\src}[1]{{\tiny\color{black!55}#1}}

% ---------------------------------------------------------------------------
%% Course map: the recurring "you are here" slide for ECON 282E.
%% Loaded by course_preamble.tex. Pattern follows AI-ECON-2026 lec01_map.tex, but the map
%% shows the whole ten-session course rather than one lecture.
%%
%%   \CourseMap{s}{label}   s = 1..10 highlights session s and prints "you are here: label"
%%                          (e.g. \CourseMap{1}{1.A2}); earlier sessions get a check mark.
%%   \CourseMap{0}{}        full overview, nothing highlighted.
%%
%% Note: TikZ option lists are comma-split before TeX conditionals run, so every \ifnum
%% below sits outside [...] option lists.

\newcommand{\CourseMap}[2]{%
\begin{center}
\begin{tikzpicture}[
    sess/.style={rectangle, rounded corners=2pt, draw=black!45, minimum width=1.34cm,
        minimum height=1.12cm, text width=1.26cm, align=center, inner sep=1pt},
    stage/.style={font=\tiny\bfseries, text=black!70},
]
  \foreach \i/\d/\t/\c in {%
      1/Sep 24/Intro \\ Growth/StageTrunk,
      2/Oct 1/HA $\cdot$ OLG \\ Policy/StageBranch,
      3/Oct 8/Num.\ DP \\ Python/StageMachine,
      4/Oct 15/Numerics \\ I/StageMachine,
      5/Oct 22/Numerics \\ II/StageMachine,
      6/Oct 29/The wall \\ \textbf{Midterm}/StageWall,
      7/Nov 5/ML \\ Deep nets/StageTools,
      8/Nov 12/RL \\ HA $+$ DL/StageTools,
      9/Nov 19/Text \\ LLMs/StageData,
      10/Dec 3/Agents \\ \textbf{Projects}/StageAgents} {
    \node[sess, fill=\c] (n\i) at ({1.46*(\i-1)}, 0)
      {{\scriptsize\bfseries S\i}\\[-1pt]{\tiny\color{black!60}\d}\\[1pt]{\tiny \t}};
    \ifnum\i<#1
      \node[font=\scriptsize, text=green!45!black] at ([xshift=-2.5mm,yshift=-2.2mm]n\i.north east) {$\checkmark$};
    \fi
    \ifnum\i=#1
      \draw[RedTitle, very thick, rounded corners=2pt]
        ([xshift=-0.6mm,yshift=-0.6mm]n\i.south west) rectangle ([xshift=0.6mm,yshift=0.6mm]n\i.north east);
      % keep the label inside the picture at both ends of the row
      \ifnum\i<3
        \node[font=\scriptsize\bfseries, text=RedTitle, anchor=south west] at ([yshift=1.2mm]n\i.north west)
          {$\blacktriangledown$\,you are here: #2};
      \else\ifnum\i>8
        \node[font=\scriptsize\bfseries, text=RedTitle, anchor=south east] at ([yshift=1.2mm]n\i.north east)
          {you are here: #2\,$\blacktriangledown$};
      \else
        \node[font=\scriptsize\bfseries, text=RedTitle, anchor=south] at ([yshift=1.2mm]n\i.north)
          {you are here: #2\,$\blacktriangledown$};
      \fi\fi
    \fi
  }
  % stage band under the sessions
  \foreach \a/\b/\lab/\c in {1/1/trunk/StageTrunk, 2/2/branches/StageBranch,
      3/5/the machine $\cdot$ classical methods/StageMachine, 6/6/the wall/StageWall,
      7/8/new tools/StageTools, 9/9/new data/StageData, 10/10/colleagues/StageAgents} {
    \fill[\c] ([yshift=-1.5mm]n\a.south west) rectangle ([yshift=-4.6mm]n\b.south east);
    \path ([yshift=-1.5mm]n\a.south west) -- ([yshift=-4.6mm]n\b.south east)
      node[midway, stage] {\lab};
  }
  \node[font=\footnotesize\itshape, text=black!70] at ([yshift=-9.5mm]$(n1.south)!0.5!(n10.south)$)
    {Build it on paper $\;\to\;$ solve it on the machine $\;\to\;$ push it to the frontier with AI};
\end{tikzpicture}
\end{center}
}


%%   \decktitle{1}{A1}{Who I Am \& Why This Course}{September 24, 2026}
%%   \begin{document}
%%   \makedecktitle
%%   ...
%%
%% Adapted from Courses/AI-ECON-2026/source/shared_preamble.tex (Metropolis theme,
%% concept/caution/example/takeaway boxes, \sectiondivider). Changes for this course:
%%   * course title block (\decktitle / \makedecktitle) instead of \makebeamertitle
%%   * \zh{...} typesets Chinese (book titles) under pdflatex via CJKutf8
%%   * researchbox for the "Research in the wild" frames
%%   * section pages off; TikZ libraries and the course map loaded here
%% Build with pdflatex only (two passes); tools/build_slides.py does this for you.

\documentclass[10pt,english,aspectratio=169]{beamer}

\usepackage[T1]{fontenc}
\usepackage{textcomp}
\usepackage[utf8]{inputenc}
\usepackage{babel}
\usepackage{CJKutf8}

\usepackage{amsmath, amssymb, amsthm, graphicx}
\usepackage{booktabs, multirow, tabularx}
\usepackage{tikz}
\usetikzlibrary{positioning,arrows.meta,shapes,calc,fit,backgrounds}
\usepackage{pgfplots}
\pgfplotsset{compat=1.16}
\usepackage[most]{tcolorbox}
\tcbuselibrary{skins, breakable}
\usepackage{listings}
\usepackage{xcolor}

\lstset{
  basicstyle=\ttfamily\small,
  keywordstyle=\color{blue!80!black}\bfseries,
  commentstyle=\color{green!50!black}\itshape,
  stringstyle=\color{orange!80!black},
  backgroundcolor=\color{gray!8},
  frame=single,
  rulecolor=\color{gray!40},
  breaklines=true,
  showstringspaces=false,
  numbers=none,
}

\usetheme{metropolis}
\usefonttheme{professionalfonts}
\metroset{sectionpage=none}

\definecolor{LightGray}{RGB}{230,230,230}
\definecolor{RedTitle}{RGB}{0,0,200}
\definecolor{VeryLightGray}{RGB}{245,245,245}
\definecolor{DarkText}{RGB}{0,0,0}
\definecolor{UCRBlue}{RGB}{0,45,106}
\definecolor{UCRGold}{RGB}{241,171,0}

% Stage colours of the course arc (used by the course map and elsewhere)
\colorlet{StageTrunk}{blue!12}
\colorlet{StageBranch}{cyan!14}
\colorlet{StageMachine}{gray!18}
\colorlet{StageWall}{red!14}
\colorlet{StageTools}{orange!20}
\colorlet{StageData}{green!16}
\colorlet{StageAgents}{violet!14}

\hypersetup{bookmarks=false,breaklinks=true,pdfborder={0 0 0},colorlinks=true,
  linkcolor=DarkText,urlcolor=RedTitle}

\setbeamercolor{background canvas}{bg=white}
\setbeamercolor{title}{fg=RedTitle}
\setbeamercolor{frametitle}{bg=VeryLightGray,fg=DarkText}
\setbeamercolor{normal text}{fg=DarkText}
\setbeamercolor{alerted text}{fg=RedTitle}
\setbeamersize{text margin left=5mm, text margin right=5mm}
\addtobeamertemplate{frametitle}{}{\vspace{-0.5em}}

\setbeamertemplate{navigation symbols}{}
\setbeamertemplate{footline}{%
  \hspace{5mm}{\usebeamerfont{page number in head/foot}\color{black!45}%
  ECON 282E \,$\cdot$\, \insertshortdeck}\hfill%
  \usebeamercolor[fg]{page number in head/foot}%
  \usebeamerfont{page number in head/foot}%
  \insertframenumber\,/\,\inserttotalframenumber\kern1em\vskip2pt%
}

% ---------------------------------------------------------------------------
% Chinese text under pdflatex (book titles etc.)
\newcommand{\zh}[1]{\begin{CJK*}{UTF8}{gbsn}#1\end{CJK*}}

% ---------------------------------------------------------------------------
% Deck title block
%   \decktitle{session}{deck id}{title}{date}
\newcommand{\insertshortdeck}{}
\newcommand{\decksession}{}
\newcommand{\deckid}{}
\newcommand{\decktitletext}{}
\newcommand{\deckdate}{}
\newcommand{\decktitle}[4]{%
  \renewcommand{\decksession}{#1}%
  \renewcommand{\deckid}{#2}%
  \renewcommand{\decktitletext}{#3}%
  \renewcommand{\deckdate}{#4}%
  \renewcommand{\insertshortdeck}{Session #1 \,$\cdot$\, Deck #1.#2}%
  \title{#3}%
  \author{Zhigang Feng}%
  \date{#4}%
}
\newcommand{\makedecktitle}{%
  \begin{frame}[plain,noframenumbering]
    \begin{tikzpicture}[remember picture,overlay]
      \fill[UCRBlue] (current page.north west) rectangle ([yshift=-0.55cm]current page.north east);
      \fill[UCRGold] ([yshift=-0.55cm]current page.north west) rectangle ([yshift=-0.62cm]current page.north east);
      \node[anchor=west, text=white, font=\small\bfseries] at ([xshift=5mm,yshift=-0.275cm]current page.north west)
        {ECON 282E \,$\cdot$\, Foundations of Macroeconomics};
      \node[anchor=east, text=white, font=\small] at ([xshift=-5mm,yshift=-0.275cm]current page.north east)
        {UC Riverside \,$\cdot$\, Fall 2026};
    \end{tikzpicture}
    \vspace*{1.1cm}
    {\usebeamerfont{title}\color{black!55}\large Session \decksession\ \,$\cdot$\, Deck \decksession.\deckid\par}
    \vspace{0.25cm}
    {\usebeamerfont{title}\color{RedTitle}\LARGE\bfseries \decktitletext\par}
    \vspace{0.7cm}
    {\large Zhigang Feng\par}
    \vspace{0.15cm}
    {\color{black!65}\deckdate\par}
    \vfill
    {\footnotesize\itshape\color{black!55}Build it on paper $\;\to\;$ solve it on the machine $\;\to\;$ push it to the frontier with AI\par}
  \end{frame}}

% ---------------------------------------------------------------------------
% Section divider (plain frame, not counted)
\newcommand{\sectiondivider}[2]{%
\begin{frame}[plain,noframenumbering]
\begin{center}
\vfill
{\Large\textbf{#1}\par}
\vspace{0.35cm}
{\normalsize #2\par}
\vfill
\end{center}
\end{frame}
}

% ---------------------------------------------------------------------------
% Boxes
\newtcolorbox{conceptbox}[1][]{colback=blue!3, colframe=RedTitle!55, boxrule=0.35mm,
  arc=1mm, left=2mm, right=2mm, top=1mm, bottom=1mm, #1}
\newtcolorbox{cautionbox}[1][]{colback=red!3, colframe=red!50!black, boxrule=0.35mm,
  arc=1mm, left=2mm, right=2mm, top=1mm, bottom=1mm, #1}
\newtcolorbox{examplebox}[1][]{colback=green!3, colframe=green!45!black, boxrule=0.35mm,
  arc=1mm, left=2mm, right=2mm, top=1mm, bottom=1mm, #1}
\newtcolorbox{takeawaybox}[1][]{colback=VeryLightGray, colframe=RedTitle!55, boxrule=0.35mm,
  arc=1mm, left=2mm, right=2mm, top=1mm, bottom=1mm, #1}
\newtcolorbox{researchbox}[1][]{enhanced, colback=violet!4, colframe=violet!60!black,
  boxrule=0.35mm, arc=1mm, left=2mm, right=2mm, top=1mm, bottom=1mm,
  fonttitle=\bfseries\small, title={Research in the wild}, #1}

\newcommand{\compacttables}{\renewcommand{\arraystretch}{1.15}}
\newcommand{\src}[1]{{\tiny\color{black!55}#1}}

% ---------------------------------------------------------------------------
%% Course map: the recurring "you are here" slide for ECON 282E.
%% Loaded by course_preamble.tex. Pattern follows AI-ECON-2026 lec01_map.tex, but the map
%% shows the whole ten-session course rather than one lecture.
%%
%%   \CourseMap{s}{label}   s = 1..10 highlights session s and prints "you are here: label"
%%                          (e.g. \CourseMap{1}{1.A2}); earlier sessions get a check mark.
%%   \CourseMap{0}{}        full overview, nothing highlighted.
%%
%% Note: TikZ option lists are comma-split before TeX conditionals run, so every \ifnum
%% below sits outside [...] option lists.

\newcommand{\CourseMap}[2]{%
\begin{center}
\begin{tikzpicture}[
    sess/.style={rectangle, rounded corners=2pt, draw=black!45, minimum width=1.34cm,
        minimum height=1.12cm, text width=1.26cm, align=center, inner sep=1pt},
    stage/.style={font=\tiny\bfseries, text=black!70},
]
  \foreach \i/\d/\t/\c in {%
      1/Sep 24/Intro \\ Growth/StageTrunk,
      2/Oct 1/HA $\cdot$ OLG \\ Policy/StageBranch,
      3/Oct 8/Num.\ DP \\ Python/StageMachine,
      4/Oct 15/Numerics \\ I/StageMachine,
      5/Oct 22/Numerics \\ II/StageMachine,
      6/Oct 29/The wall \\ \textbf{Midterm}/StageWall,
      7/Nov 5/ML \\ Deep nets/StageTools,
      8/Nov 12/RL \\ HA $+$ DL/StageTools,
      9/Nov 19/Text \\ LLMs/StageData,
      10/Dec 3/Agents \\ \textbf{Projects}/StageAgents} {
    \node[sess, fill=\c] (n\i) at ({1.46*(\i-1)}, 0)
      {{\scriptsize\bfseries S\i}\\[-1pt]{\tiny\color{black!60}\d}\\[1pt]{\tiny \t}};
    \ifnum\i<#1
      \node[font=\scriptsize, text=green!45!black] at ([xshift=-2.5mm,yshift=-2.2mm]n\i.north east) {$\checkmark$};
    \fi
    \ifnum\i=#1
      \draw[RedTitle, very thick, rounded corners=2pt]
        ([xshift=-0.6mm,yshift=-0.6mm]n\i.south west) rectangle ([xshift=0.6mm,yshift=0.6mm]n\i.north east);
      % keep the label inside the picture at both ends of the row
      \ifnum\i<3
        \node[font=\scriptsize\bfseries, text=RedTitle, anchor=south west] at ([yshift=1.2mm]n\i.north west)
          {$\blacktriangledown$\,you are here: #2};
      \else\ifnum\i>8
        \node[font=\scriptsize\bfseries, text=RedTitle, anchor=south east] at ([yshift=1.2mm]n\i.north east)
          {you are here: #2\,$\blacktriangledown$};
      \else
        \node[font=\scriptsize\bfseries, text=RedTitle, anchor=south] at ([yshift=1.2mm]n\i.north)
          {you are here: #2\,$\blacktriangledown$};
      \fi\fi
    \fi
  }
  % stage band under the sessions
  \foreach \a/\b/\lab/\c in {1/1/trunk/StageTrunk, 2/2/branches/StageBranch,
      3/5/the machine $\cdot$ classical methods/StageMachine, 6/6/the wall/StageWall,
      7/8/new tools/StageTools, 9/9/new data/StageData, 10/10/colleagues/StageAgents} {
    \fill[\c] ([yshift=-1.5mm]n\a.south west) rectangle ([yshift=-4.6mm]n\b.south east);
    \path ([yshift=-1.5mm]n\a.south west) -- ([yshift=-4.6mm]n\b.south east)
      node[midway, stage] {\lab};
  }
  \node[font=\footnotesize\itshape, text=black!70] at ([yshift=-9.5mm]$(n1.south)!0.5!(n10.south)$)
    {Build it on paper $\;\to\;$ solve it on the machine $\;\to\;$ push it to the frontier with AI};
\end{tikzpicture}
\end{center}
}


\graphicspath{{./figures/}}

\lstset{language=Python, basicstyle=\ttfamily\scriptsize, frame=none,
        backgroundcolor=\color{gray!8}, aboveskip=2pt, belowskip=2pt}

\decktitle{5}{B5}{Determining the Coefficients}{Thursday, October 22, 2026}

\begin{document}

\makedecktitle

%=============================================================================
\begin{frame}{Where We Are}
\CourseMap{5}{5.B5}
\vspace{-1mm}
\begin{center}
\small \textbf{Block B:} 5.B1 the residual $\;\cdot\;$ 5.B2 choosing a basis $\;\cdot\;$ 5.B3 Chebyshev $\;\cdot\;$ 5.B4 finite elements $\;\cdot\;$ \textbf{5.B5} determining the coefficients $\;\cdot\;$ 5.B6 the worked example.
\end{center}
\end{frame}

%=============================================================================
\begin{frame}{Weighted Residuals}
\begin{columns}[T]
\begin{column}{0.55\textwidth}
\small
The standard answer to ``what does small mean'' is: \textbf{zero in a weighted integral sense}. Choose $n+1$ weight functions $w_i:\Omega\to\mathbb{R}^{m}$ and require
\[
  \int_{\Omega} w_i(x)\,R(x;\theta)\,dx \;=\; 0,
  \qquad i=0,\dots,n .
\]
That is $n+1$ equations for $n+1$ coefficients --- square, and generically solvable.
\medskip

This is why projection methods are also called \textbf{weighted residual methods}. Every named method in the literature is one choice of $\{w_i\}$.
\end{column}
\begin{column}{0.42\textwidth}
\begin{takeawaybox}\footnotesize
And the payoff: an \textbf{intractable functional equation} has become a \textbf{standard nonlinear system}. Newton for small problems, conjugate gradient for large ones --- all of 4.A1 and 4.A3, unchanged.
\end{takeawaybox}
\medskip
\begin{examplebox}\footnotesize
Notation: the literature writes the weights as $f_i$ and the basis as $\psi_i$, $g_i$ or $\varphi_i$ depending on the author. Here the basis is always $\psi_i$ and the weights always $w_i$.
\end{examplebox}
\end{column}
\end{columns}
\vspace{1mm}
\src{Quant\_Macro Lec.~9, ``Objective function''.}
\end{frame}

%=============================================================================
\begin{frame}{What the Weight Is Actually For}
\begin{columns}[T]
\begin{column}{0.55\textwidth}
\small
A policy function does not misbehave uniformly. It may bend sharply near a constraint and be almost linear elsewhere --- so a residual of a given size means different things in different regions.
\medskip

The weight lets you say where accuracy matters: put more weight where the function varies rapidly, and the solve will spend its coefficients there.
\medskip

The source's analogy is apt: the weight handles \textbf{``heteroskedasticity'' in the function being approximated}. It is the same instinct as weighted least squares, and the same instinct as 3.A2's log-spaced asset grid.
\end{column}
\begin{column}{0.42\textwidth}
\begin{conceptbox}[title=\textbf{Basis functions as weights}]\small
A natural choice is to use the \textbf{basis itself}. If the basis was chosen to match the features of the function, it already puts its variation where the function's is --- so weighting by it emphasises exactly the right regions.
\medskip
That choice has a name: \textbf{Galerkin}.
\end{conceptbox}
\end{column}
\end{columns}
\vspace{1mm}
\src{Quant\_Macro Lec.~9, ``Weight function'' (both frames).}
\end{frame}

%=============================================================================
\begin{frame}{I. Least Squares}
\begin{columns}[T]
\begin{column}{0.55\textwidth}
\small
Take $\;w_i(x)=\dfrac{\partial R(x;\theta)}{\partial\theta_i}$.
\medskip

This is not arbitrary --- it is the first-order condition of the variational problem
\[
  \min_{\theta}\int_{\Omega}R^{2}(x;\theta)\,dx,
\]
whose FOC is exactly $\int_\Omega \partial_{\theta_i}R\cdot R\,dx=0$.
\medskip

So least squares is the \textbf{regression} answer: minimise the squared residual over the domain, and you are choosing the weights without realising it.
\end{column}
\begin{column}{0.42\textwidth}
\begin{conceptbox}[title=\textbf{Why it is attractive}]\small
It always generates \textbf{symmetric} matrices --- convenient in the proofs, and exploitable by algorithms that are faster and use less memory.
\end{conceptbox}
\medskip
\begin{cautionbox}\footnotesize
And why it is not: least squares is prone to \textbf{ill-conditioning}, and the resulting systems can be awkward to solve. Squaring the residual squares the condition number too.
\end{cautionbox}
\end{column}
\end{columns}
\vspace{1mm}
\src{Quant\_Macro Lec.~9, ``Weight function I: Least Squares''.}
\end{frame}

%=============================================================================
\begin{frame}{II. Subdomain, and III. Moments}
\begin{columns}[T]
\begin{column}{0.49\textwidth}
\small
\textbf{Subdomain.} Split $\Omega$ into $n+1$ pieces $\Omega_i$ and take step functions,
\[
  w_i(x)=\mathbf{1}\{x\in\Omega_i\},
\]
so the condition is that the residual \textbf{averages to zero on each piece}. It may be large inside a piece, as long as it cancels.
\medskip

\textbf{Moments.} Take $\{1,x,x^{2},\dots,x^{n}\}$ as weights:
\[
  \int_{\Omega}x^{i}R(x;\theta)\,dx=0 .
\]
The first $n$ moments of the residual are zero.
\end{column}
\begin{column}{0.48\textwidth}
\begin{cautionbox}[title=\textbf{Read it carefully}]\footnotesize
The weights are $\{1,x,x^{2},\dots\}$ --- the first one is the constant, so the first condition just says the residual \emph{integrates to zero}. The rest set its higher moments to zero.
\end{cautionbox}
\medskip
\begin{examplebox}\footnotesize
The verdict on moments is 5.B2's argument again: fine for $n=2$ or $3$, and for higher orders the monomials are so collinear that rounding error dominates. \textbf{``Moments are to be avoided.''}
\end{examplebox}
\end{column}
\end{columns}
\vspace{1mm}
\src{Quant\_Macro Lec.~9, ``Weight Function II: Subdomain Approach'' and ``Weight function III: Moments''.}
\end{frame}

%=============================================================================
\begin{frame}{IV. Collocation}
\begin{columns}[T]
\begin{column}{0.55\textwidth}
\small
Take the weights to be Dirac deltas:
\[
  w_i(x)=\delta(x-x_i).
\]
Since $\int \delta(x-x_i)R(x)\,dx = R(x_i)$, the conditions collapse to
\[
  R(x_i;\theta)=0,\qquad i=0,\dots,n .
\]
\textbf{The residual is exactly zero at $n+1$ points and unconstrained in between} --- 5.B1's picture.
\medskip

The attraction is cost: no integral is ever computed. Each condition is one function evaluation, which matters when $R$ is expensive and strongly nonlinear.
\end{column}
\begin{column}{0.42\textwidth}
\begin{conceptbox}[title=\textbf{Picking the points}]\footnotesize
A systematic choice uses a density $\mu_\gamma$, placing $x_j$ so that $\int_{-1}^{x_j}\mu_\gamma=j/n$. Larger $\gamma$ clusters points at the endpoints; equi-spacing is the limit $\gamma\to0$.
\end{conceptbox}
\medskip
\begin{cautionbox}\footnotesize
Read the limit carefully: $\gamma=0$ substituted into $\mu_\gamma$ gives \textbf{the zero function}, which is not a density. Equi-spacing is the limit $\gamma\to0$, not the value $\gamma=0$.
\end{cautionbox}
\end{column}
\end{columns}
\vspace{1mm}
\src{Quant\_Macro Lec.~9, ``Weight function IV: Collocation or pseudospectral'' and ``Intuition: $\mu_\gamma(x)$ and the Case $\gamma=0$''.}
\end{frame}

%=============================================================================
\begin{frame}{Collocation Against Subdomain}
\small
\begin{center}
\begin{tabularx}{0.95\linewidth}{@{}>{\raggedright\arraybackslash}p{2.0cm} >{\raggedright\arraybackslash}X >{\raggedright\arraybackslash}X@{}}
\toprule
 & \textbf{Collocation} & \textbf{Subdomain} \\
\midrule
Weight & $\delta(x-x_i)$ & step function on $\Omega_i$ \\
Residual is & zero \textbf{at points} & zero \textbf{on average} over each piece \\
Evaluation & pointwise only & an integral per condition \\
Cost & \textbf{lower} & higher \\
Best for & strongly \textbf{nonlinear}, spectral bases & \textbf{smooth}, PDE-like problems \\
Limitation & says nothing between the points & integration overhead \\
\bottomrule
\end{tabularx}
\end{center}
\medskip
\begin{center}
\begin{minipage}{0.88\linewidth}
\begin{conceptbox}[title=\textbf{The shared weakness}]\footnotesize
Neither controls the residual \emph{between} the conditions. Collocation leaves it free between points; subdomain lets it cancel within a piece. That is why you always grade the finished solution on a \textbf{different, finer set of points} --- 3.A1's separation of grids.
\end{conceptbox}
\end{minipage}
\end{center}
\vspace{1mm}
\src{Quant\_Macro Lec.~9, ``Dirac Delta Function $\delta$ and Comparison''.}
\end{frame}

%=============================================================================
\begin{frame}{V. Orthogonal Collocation, and VI. Galerkin}
\begin{columns}[T]
\begin{column}{0.5\textwidth}
\small
\textbf{Orthogonal collocation.} Collocation, with two extra commitments: the basis is a family of \textbf{orthogonal polynomials}, and the collocation points are the \textbf{roots of the $n$-th one}.
\medskip

Those roots are not equi-spaced --- they cluster at the endpoints, which is exactly where 5.B3 said an equi-spaced rule fails. The result converges exponentially for smooth functions, and the standing verdict in this literature is ``surprisingly good performance''.
\medskip

\textbf{Galerkin.} Make the residual \textbf{orthogonal to every basis function}:
\[
  \int_{\Omega}\psi_i(x)\,R(x;\theta)\,dx=0 .
\]
\end{column}
\begin{column}{0.47\textwidth}
\begin{conceptbox}[title=\textbf{What Galerkin means}]\small
The residual has \textbf{no component the basis could still represent}. Anything left over is invisible in this function space --- so you have extracted everything the basis can give.
\end{conceptbox}
\medskip
\begin{takeawaybox}\footnotesize
If the basis is complete, the Galerkin solution converges \textbf{pointwise} to the truth. It is accurate and robust, and \emph{difficult to code}: every condition is an integral of a nonlinear function.
\medskip
Rule of thumb: Galerkin at order $n$ is about as accurate as collocation at $n+1$ or $n+2$.
\end{takeawaybox}
\end{column}
\end{columns}
\vspace{1mm}
\src{Quant\_Macro Lec.~9, ``Weight function V: Orthogonal collocation'' and ``VI: Galerkin or Rayleigh-Ritz''.}
\end{frame}

%=============================================================================
\begin{frame}{Measured: Six Weight Functions, One Model, One Basis}
\begin{columns}[T]
\begin{column}{0.56\textwidth}
\small
The quarterly RBC of Block A, Chebyshev degree $6$ in $k$ for each of $7$ shock states --- \textbf{49 unknowns}, identical in every row. Only the weight function changes. All six converge; integrals by $40$-node Gauss--Legendre.
\medskip
\begin{center}\footnotesize
\begin{tabularx}{\linewidth}{@{}>{\raggedright\arraybackslash}X r r@{}}
\toprule
\textbf{Weight} & $\max|\mathcal{E}|$ & \textbf{sec} \\
\midrule
orthogonal collocation & $\mathbf{10^{-5.78}}$ & $0.11$ \\
least squares & $10^{-5.61}$ & $0.80$ \\
Galerkin & $10^{-5.55}$ & $0.20$ \\
moments & $10^{-5.55}$ & $0.15$ \\
collocation, equi-spaced & $10^{-5.39}$ & $0.11$ \\
subdomain & $10^{-5.05}$ & $0.14$ \\
\bottomrule
\end{tabularx}
\end{center}
\end{column}
\begin{column}{0.41\textwidth}
\begin{takeawaybox}\footnotesize
\textbf{Every method lands within $0.73$ of a decade of every other} --- a factor of $5.4$ between
best and worst. The choice of weight function is worth less than one extra basis coefficient, which
5.B6 measures at $1.2$ decades for two.
\end{takeawaybox}
\medskip
\begin{examplebox}\footnotesize
And the cheapest is the best: orthogonal collocation wins on accuracy \emph{and} ties for fastest, because it evaluates no integrals at all.
\end{examplebox}
\end{column}
\end{columns}
\vspace{1mm}
\src{Ledger \texttt{collocation}. Warm-started from the first-order perturbation solution; residual norms $\le 1.4\times10^{-14}$ except least squares, $6.9\times10^{-7}$.}
\end{frame}

%=============================================================================
\begin{frame}{Two Things the Table Says}
\begin{columns}[T]
\begin{column}{0.54\textwidth}
\small
\textbf{1. Moments and Galerkin are the same method here.} Their errors agree to nine decimal places: $-5.5515379662$ against $-5.5515379670$.
\medskip

That is not luck. $\{1,x,\dots,x^{n}\}$ and $\{T_0,\dots,T_n\}$ \textbf{span the same space}, so ``orthogonal to every monomial'' and ``orthogonal to every Chebyshev polynomial'' are the same $n+1$ linear conditions. In exact arithmetic they are identical.
\medskip

The warning against moments is therefore \textbf{purely numerical}: the conditions are right, and the basis you write them in destroys them. At $n=6$ that has not bitten yet; 5.B2's Hilbert table says where it will.
\end{column}
\begin{column}{0.43\textwidth}
\begin{cautionbox}[title=\textbf{2. Least squares is 7x slower}]\small
$0.80$\,s against $0.11$\,s, for no accuracy gain --- and it is the one method that did not drive its residual to machine zero, stopping at $6.9\times10^{-7}$.
\medskip
Exactly the predicted ill-conditioning, showing up as a solver that will not converge the last six digits.
\end{cautionbox}
\medskip
\begin{examplebox}\footnotesize
The source's ``Galerkin $\approx$ collocation at $n+1$ or $n+2$'' is not what we measure: here \emph{orthogonal} collocation beats Galerkin at the same $n$.
\end{examplebox}
\end{column}
\end{columns}
\vspace{1mm}
\src{Ledger \texttt{collocation.rows}; full precision in the JSON.}
\end{frame}

%=============================================================================
\begin{frame}{Takeaways}
\begin{takeawaybox}
\begin{itemize}\small
  \item ``Make the residual small'' becomes $\int_\Omega w_i R\,dx=0$ for $n+1$ weights --- a square \textbf{nonlinear system}, solved with the tools of 4.A1 and 4.A3. Every named projection method is one choice of $\{w_i\}$.
  \item \textbf{Least squares} is the regression answer and gives symmetric matrices; it is also the most ill-conditioned, and was measured $7\times$ slower with the worst final residual.
  \item \textbf{Subdomain} zeroes the residual on average over pieces; \textbf{moments} zeroes its first $n$ moments and is to be avoided at high order; \textbf{collocation} zeroes it at points and computes no integrals at all.
  \item \textbf{Galerkin} makes the residual orthogonal to the basis --- nothing is left that the basis could represent --- and converges pointwise if the basis is complete.
  \item Measured on one model with one basis: \textbf{all six land within $0.73$ of a decade}. The weight function matters less than one extra coefficient.
  \item \textbf{Orthogonal collocation wins on accuracy and ties for fastest}, which is why it is the default in this literature.
  \item And \textbf{moments and Galerkin agreed to nine decimals} --- they are the same conditions in two bases. The objection to moments is conditioning, not mathematics.
\end{itemize}
\end{takeawaybox}
\end{frame}

%=============================================================================
\begin{frame}{Bridge: Now Solve the Model}
\begin{columns}[T]
\begin{column}{0.53\textwidth}
\small
Block B now has both halves: a basis (Chebyshev, on a domain, with clustered nodes) and a projection (orthogonal collocation, because it is cheapest and best).
\medskip

That is a complete method. The last deck runs it on the model Sessions 3, 4 and 5 have all solved --- and then puts every method in this course into one table.
\end{column}
\begin{column}{0.44\textwidth}
\begin{conceptbox}[title=\textbf{Next (5.B6)}]\small
The worked example: the same quarterly RBC by Chebyshev collocation, five steps, warm-started from perturbation. Then \textbf{one model, four methods} --- and the answer to which one you should reach for.
\end{conceptbox}
\end{column}
\end{columns}
\end{frame}

\end{document}
